% Options for packages loaded elsewhere
% Options for packages loaded elsewhere
\PassOptionsToPackage{unicode}{hyperref}
\PassOptionsToPackage{hyphens}{url}
%
\documentclass[
  english,
  russian,
  12pt,
  a4paper,
  DIV=11,
  numbers=noendperiod]{scrreprt}
\usepackage{xcolor}
\usepackage{amsmath,amssymb}
\setcounter{secnumdepth}{5}
\usepackage{iftex}
\ifPDFTeX
  \usepackage[T1]{fontenc}
  \usepackage[utf8]{inputenc}
  \usepackage{textcomp} % provide euro and other symbols
\else % if luatex or xetex
  \usepackage{unicode-math} % this also loads fontspec
  \defaultfontfeatures{Scale=MatchLowercase}
  \defaultfontfeatures[\rmfamily]{Ligatures=TeX,Scale=1}
\fi
\usepackage{lmodern}
\ifPDFTeX\else
  % xetex/luatex font selection
\fi
% Use upquote if available, for straight quotes in verbatim environments
\IfFileExists{upquote.sty}{\usepackage{upquote}}{}
\IfFileExists{microtype.sty}{% use microtype if available
  \usepackage[]{microtype}
  \UseMicrotypeSet[protrusion]{basicmath} % disable protrusion for tt fonts
}{}
\usepackage{setspace}
% Make \paragraph and \subparagraph free-standing
\makeatletter
\ifx\paragraph\undefined\else
  \let\oldparagraph\paragraph
  \renewcommand{\paragraph}{
    \@ifstar
      \xxxParagraphStar
      \xxxParagraphNoStar
  }
  \newcommand{\xxxParagraphStar}[1]{\oldparagraph*{#1}\mbox{}}
  \newcommand{\xxxParagraphNoStar}[1]{\oldparagraph{#1}\mbox{}}
\fi
\ifx\subparagraph\undefined\else
  \let\oldsubparagraph\subparagraph
  \renewcommand{\subparagraph}{
    \@ifstar
      \xxxSubParagraphStar
      \xxxSubParagraphNoStar
  }
  \newcommand{\xxxSubParagraphStar}[1]{\oldsubparagraph*{#1}\mbox{}}
  \newcommand{\xxxSubParagraphNoStar}[1]{\oldsubparagraph{#1}\mbox{}}
\fi
\makeatother


\usepackage{longtable,booktabs,array}
\usepackage{calc} % for calculating minipage widths
% Correct order of tables after \paragraph or \subparagraph
\usepackage{etoolbox}
\makeatletter
\patchcmd\longtable{\par}{\if@noskipsec\mbox{}\fi\par}{}{}
\makeatother
% Allow footnotes in longtable head/foot
\IfFileExists{footnotehyper.sty}{\usepackage{footnotehyper}}{\usepackage{footnote}}
\makesavenoteenv{longtable}
\usepackage{graphicx}
\makeatletter
\newsavebox\pandoc@box
\newcommand*\pandocbounded[1]{% scales image to fit in text height/width
  \sbox\pandoc@box{#1}%
  \Gscale@div\@tempa{\textheight}{\dimexpr\ht\pandoc@box+\dp\pandoc@box\relax}%
  \Gscale@div\@tempb{\linewidth}{\wd\pandoc@box}%
  \ifdim\@tempb\p@<\@tempa\p@\let\@tempa\@tempb\fi% select the smaller of both
  \ifdim\@tempa\p@<\p@\scalebox{\@tempa}{\usebox\pandoc@box}%
  \else\usebox{\pandoc@box}%
  \fi%
}
% Set default figure placement to htbp
\def\fps@figure{htbp}
\makeatother



\ifLuaTeX
\usepackage[bidi=basic,provide=*]{babel}
\else
\usepackage[bidi=default,provide=*]{babel}
\fi
% get rid of language-specific shorthands (see #6817):
\let\LanguageShortHands\languageshorthands
\def\languageshorthands#1{}


\setlength{\emergencystretch}{3em} % prevent overfull lines

\providecommand{\tightlist}{%
  \setlength{\itemsep}{0pt}\setlength{\parskip}{0pt}}



 
\usepackage[style=gost-numeric,backend=biber,langhook=extras,autolang=other*]{biblatex}
\addbibresource{bib/cite.bib}

\usepackage[]{csquotes}

\usepackage{indentfirst}
\usepackage{float}
\floatplacement{figure}{H}
\IfFileExists{plex-otf.sty}{
  %% Full TeXlive
  \usepackage[
    % math,
    RM={Scale=0.94},SS={Scale=0.94},SScon={Scale=0.94},TT={Scale=MatchLowercase,FakeStretch=0.9},
    DefaultFeatures={Ligatures=Common}
  ]{plex-otf}
  % \usepackage[Scale=MatchUppercase]{juliamono}
}{
  %% TinyTeX
  \usepackage{libertine}
}
\KOMAoption{captions}{tableheading}
\makeatletter
\@ifpackageloaded{caption}{}{\usepackage{caption}}
\AtBeginDocument{%
\ifdefined\contentsname
  \renewcommand*\contentsname{Содержание}
\else
  \newcommand\contentsname{Содержание}
\fi
\ifdefined\listfigurename
  \renewcommand*\listfigurename{Список иллюстраций}
\else
  \newcommand\listfigurename{Список иллюстраций}
\fi
\ifdefined\listtablename
  \renewcommand*\listtablename{Список таблиц}
\else
  \newcommand\listtablename{Список таблиц}
\fi
\ifdefined\figurename
  \renewcommand*\figurename{Рисунок}
\else
  \newcommand\figurename{Рисунок}
\fi
\ifdefined\tablename
  \renewcommand*\tablename{Таблица}
\else
  \newcommand\tablename{Таблица}
\fi
}
\@ifpackageloaded{float}{}{\usepackage{float}}
\floatstyle{ruled}
\@ifundefined{c@chapter}{\newfloat{codelisting}{h}{lop}}{\newfloat{codelisting}{h}{lop}[chapter]}
\floatname{codelisting}{Список}
\newcommand*\listoflistings{\listof{codelisting}{Листинги}}
\makeatother
\makeatletter
\makeatother
\makeatletter
\@ifpackageloaded{caption}{}{\usepackage{caption}}
\@ifpackageloaded{subcaption}{}{\usepackage{subcaption}}
\makeatother
\usepackage{bookmark}
\IfFileExists{xurl.sty}{\usepackage{xurl}}{} % add URL line breaks if available
\urlstyle{same}
\hypersetup{
  pdftitle={отчёт по лабораторной работе №12},
  pdfauthor={Кхари Жекка Кализая арсе},
  pdflang={ru-RU},
  hidelinks,
  pdfcreator={LaTeX via pandoc}}


\title{отчёт по лабораторной работе №12}
\usepackage{etoolbox}
\makeatletter
\providecommand{\subtitle}[1]{% add subtitle to \maketitle
  \apptocmd{\@title}{\par {\large #1 \par}}{}{}
}
\makeatother
\subtitle{Настройка NAT}
\author{Кхари Жекка Кализая арсе}
\date{}
\begin{document}
\maketitle

\renewcommand*\contentsname{Содержание}
{
\setcounter{tocdepth}{1}
\tableofcontents
}
\listoffigures
\listoftables

\setstretch{1.5}
\chapter{Цель
работы}\label{ux446ux435ux43bux44c-ux440ux430ux431ux43eux442ux44b}

Приобретение практических навыков по настройке доступа локальной сети к
внешней сети посредством NAT.

\chapter{Задание}\label{ux437ux430ux434ux430ux43dux438ux435}

\begin{enumerate}
\def\labelenumi{\arabic{enumi}.}
\tightlist
\item
  Сделать первоначальную настройку маршрутизатора provider-gw-1 и
  коммутатора provider-sw-1 провайдера: задать имя, настроить доступ по
  паролю и т.п. (см. разделы 12.4.1, 12.4.2).
\item
  Настроить интерфейсы маршрутизатора provider-gw-1 и коммутатора
  provider-sw-1 провайдера: (см. разделы 12.4.3, 12.4.4).
\item
  Настроить интерфейсы маршрутизатора сети «Донская» для доступа к сети
  провайдера (cм. раздел 12.4.5).
\item
  Настроить на маршрутизаторе сети «Донская» NAT с правилами, указанными
  в разделе 12.2 (см. разделы 12.4.6--12.4.8).
\item
  Настроить доступ из внешней сети в локальную сеть организации, как
  указано в разделе 12.2 (см. раздел 12.4.9).
\item
  Проверить работоспособность заданных настроек.
\item
  При выполнении работы необходимо учитывать соглашение об именовании
  (см. раздел 2.5).
\end{enumerate}

\chapter{Выполнение лабораторной
работы}\label{ux432ux44bux43fux43eux43bux43dux435ux43dux438ux435-ux43bux430ux431ux43eux440ux430ux442ux43eux440ux43dux43eux439-ux440ux430ux431ux43eux442ux44b}

сначала этой лабораторной работы я скопировал файл проекта и переназвал
его на lab12

\section{Первоначальная настройка маршрутизатора
provider-gw-1}\label{ux43fux435ux440ux432ux43eux43dux430ux447ux430ux43bux44cux43dux430ux44f-ux43dux430ux441ux442ux440ux43eux439ux43aux430-ux43cux430ux440ux448ux440ux443ux442ux438ux437ux430ux442ux43eux440ux430-provider-gw-1}

Чтобы настройть маршрутизатор я открыл его терминал и начал настройку. я
ввел следующие команды, которые установлят пароль, количество сеансов,
учетные записи шифрованный пароли (рис.~\ref{fig-002}).

\begin{figure}

\centering{

\includegraphics[width=0.7\linewidth,height=\textheight,keepaspectratio]{image/fig00002.png}

}

\caption{\label{fig-002}Первоначальная настройка маршрутизатора
provider-gw-1}

\end{figure}%

\section{Первоначальная настройка коммутатора
provider-sw-1}\label{ux43fux435ux440ux432ux43eux43dux430ux447ux430ux43bux44cux43dux430ux44f-ux43dux430ux441ux442ux440ux43eux439ux43aux430-ux43aux43eux43cux43cux443ux442ux430ux442ux43eux440ux430-provider-sw-1}

я повторил кофигурацию на коммутаторе provider-sw-1
(рис.~\ref{fig-003}).

\begin{figure}

\centering{

\includegraphics[width=0.7\linewidth,height=\textheight,keepaspectratio]{image/fig00003.png}

}

\caption{\label{fig-003}Первоначальная настройка коммутатора
provider-sw-1}

\end{figure}%

\section{Настройка интерфейсов маршрутизатора
provider-gw-1}\label{ux43dux430ux441ux442ux440ux43eux439ux43aux430-ux438ux43dux442ux435ux440ux444ux435ux439ux441ux43eux432-ux43cux430ux440ux448ux440ux443ux442ux438ux437ux430ux442ux43eux440ux430-provider-gw-1}

потом я начал настроить интерфейсы f0/0 ,f0/0.4 , f0/1.
(рис.~\ref{fig-004}).

\begin{figure}

\centering{

\includegraphics[width=0.7\linewidth,height=\textheight,keepaspectratio]{image/fig00004.png}

}

\caption{\label{fig-004}Настройка интерфейсов маршрутизатора
provider-gw-1}

\end{figure}%

\section{Настройка интерфейсов коммутатора
provider-sw-1}\label{ux43dux430ux441ux442ux440ux43eux439ux43aux430-ux438ux43dux442ux435ux440ux444ux435ux439ux441ux43eux432-ux43aux43eux43cux43cux443ux442ux430ux442ux43eux440ux430-provider-sw-1}

Дальше я настроил интерфейсы коммутатора (f0/1 , f0/1.4)
(рис.~\ref{fig-005}).

\begin{figure}

\centering{

\includegraphics[width=0.7\linewidth,height=\textheight,keepaspectratio]{image/fig00005.png}

}

\caption{\label{fig-005}Настройка интерфейсов коммутатора provider-sw-1}

\end{figure}%

\section{Настройка интерфейсов маршрутизатора
msk-donskaya-gw-1}\label{ux43dux430ux441ux442ux440ux43eux439ux43aux430-ux438ux43dux442ux435ux440ux444ux435ux439ux441ux43eux432-ux43cux430ux440ux448ux440ux443ux442ux438ux437ux430ux442ux43eux440ux430-msk-donskaya-gw-1}

Здесь я настроил интерфейсы f0/1, f0/1.4 маршрутизатора
msk-donskaya-gw-1. я включил интерфейс f0/1 потом я указать применение
стандарта IEEE 802.1Q (рис.~\ref{fig-006}).

\begin{figure}

\centering{

\includegraphics[width=0.7\linewidth,height=\textheight,keepaspectratio]{image/fig00006png.pdf}

}

\caption{\label{fig-006}Настройка интерфейсов маршрутизатора
msk-donskaya-gw-1}

\end{figure}%

\section{Настройка пула адресов для
NAT}\label{ux43dux430ux441ux442ux440ux43eux439ux43aux430-ux43fux443ux43bux430-ux430ux434ux440ux435ux441ux43eux432-ux434ux43bux44f-nat}

потом я определил группу IP-адресов, которые маршрутизатор может
использовать чтобы выполнить NAT (рис.~\ref{fig-007}).

\begin{figure}

\centering{

\includegraphics[width=0.7\linewidth,height=\textheight,keepaspectratio]{image/fig00007.png}

}

\caption{\label{fig-007}Настройка пула адресов для NAT}

\end{figure}%

\section{Настройка списка доступа для
NAT}\label{ux43dux430ux441ux442ux440ux43eux439ux43aux430-ux441ux43fux438ux441ux43aux430-ux434ux43eux441ux442ux443ux43fux430-ux434ux43bux44f-nat}

чтобы обеспечить безопасность сети, я определил список доступа для NAT
(рис.~\ref{fig-008}).

\begin{figure}

\centering{

\includegraphics[width=0.7\linewidth,height=\textheight,keepaspectratio]{image/fig00008.png}

}

\caption{\label{fig-008}Настройка списка доступа для NAT}

\end{figure}%

\section{Сеть дисплейных
классов}\label{ux441ux435ux442ux44c-ux434ux438ux441ux43fux43bux435ux439ux43dux44bux445-ux43aux43bux430ux441ux441ux43eux432}

Потом я начал настройку сети дисплейных классов (рис.~\ref{fig-009}).
для того я разрешил подключение с сети 10.128.3.0 (сеть ДК) на сервер
192.0.2.11 (www.yandex.ru) и 192.0.2.12 (stud.rudn.university)

\begin{figure}

\centering{

\includegraphics[width=0.7\linewidth,height=\textheight,keepaspectratio]{image/fig00009.png}

}

\caption{\label{fig-009}Сеть дисплейных классов}

\end{figure}%

\section{сеть
кафедр}\label{ux441ux435ux442ux44c-ux43aux430ux444ux435ux434ux440}

я повторил те же команды для сети дисплейных классов но я изменил сеть
от которого соединение выполняется (10.128.4.0 - К ) и хость к которому
идет соединение (192.0.2.13 - esystem.pfur.ru) такое соединение только
могут использовать порт 80 (рис.~\ref{fig-010}).

\begin{figure}

\centering{

\includegraphics[width=0.7\linewidth,height=\textheight,keepaspectratio]{image/fig00010.png}

}

\caption{\label{fig-010}сеть кафедр}

\end{figure}%

\section{сеть
администрации}\label{ux441ux435ux442ux44c-ux430ux434ux43cux438ux43dux438ux441ux442ux440ux430ux446ux438ux438}

Здесь разрешается подключение к серверку www.rudn.ru (192.0.2.14) с сети
администрации (рис.~\ref{fig-011}).

\begin{figure}

\centering{

\includegraphics[width=0.7\linewidth,height=\textheight,keepaspectratio]{image/fig00011.png}

}

\caption{\label{fig-011}сеть администрации}

\end{figure}%

\section{Доступ для компьютера
администратора}\label{ux434ux43eux441ux442ux443ux43f-ux434ux43bux44f-ux43aux43eux43cux43fux44cux44eux442ux435ux440ux430-ux430ux434ux43cux438ux43dux438ux441ux442ux440ux430ux442ux43eux440ux430}

Здесь я дал права доступа на все сервер компьютеру администратора
(рис.~\ref{fig-012}).

\begin{figure}

\centering{

\includegraphics[width=0.7\linewidth,height=\textheight,keepaspectratio]{image/fig00012.png}

}

\caption{\label{fig-012}Доступ для компьютера администратора}

\end{figure}%

\section{Настройка
NAT}\label{ux43dux430ux441ux442ux440ux43eux439ux43aux430-nat}

Здесь я начал настройку NAT для того я использовал команду ip nat inside
и outside чтобы определить какие сети находятся внутри и которые
находятся вне, таким образом маршрутизатор может определить какие
интерфейсы переводит (рис.~\ref{fig-014}).

\begin{figure}

\centering{

\includegraphics[width=0.7\linewidth,height=\textheight,keepaspectratio]{image/fig00014.png}

}

\caption{\label{fig-014}Настройка NAT}

\end{figure}%

\section{настройка доступа из
интернета}\label{ux43dux430ux441ux442ux440ux43eux439ux43aux430-ux434ux43eux441ux442ux443ux43fux430-ux438ux437-ux438ux43dux442ux435ux440ux43dux435ux442ux430}

здесь я настраиваю список доступа для подключений, идущих из Интернета в
подсети. для того в командах я указываю сеть, IP-адресы сервера и порты
(рис.~\ref{fig-015}).

\begin{figure}

\centering{

\includegraphics[width=0.7\linewidth,height=\textheight,keepaspectratio]{image/fig00015.png}

}

\caption{\label{fig-015}настройка доступа из интернета}

\end{figure}%

\chapter{Выводы}\label{ux432ux44bux432ux43eux434ux44b}

В этой лабораторной работе я настроил пароли в новых устройствах, также
я настроил NAT и список прав чтобы разрешить соединение с подсетей к
интернету и обратно.

\chapter*{Список
литературы}\label{ux441ux43fux438ux441ux43eux43a-ux43bux438ux442ux435ux440ux430ux442ux443ux440ux44b}
\addcontentsline{toc}{chapter}{Список литературы}

\printbibliography[heading=none]





\end{document}
